\documentclass[a4paper,12pt]{article}
\usepackage[T2A]{fontenc}
\usepackage[utf8]{inputenc}
\usepackage[russian]{babel}
\usepackage{cmap}
\usepackage{amsmath,amssymb,mathtools}
\usepackage{PTSans}
\renewcommand{\familydefault}{\sfdefault}
\usepackage{icomma}
\usepackage[a4paper,left=22mm,right=22mm,top=23mm,bottom=22mm,headheight=15pt]{geometry}
\usepackage{microtype}
\usepackage{tikz}
\usetikzlibrary{calc,arrows.meta,positioning,shapes.geometric}
\usepackage{tabularx,booktabs}
\usepackage{enumitem}
\usepackage{hyperref}
\usepackage{parskip}
\usepackage{fancyhdr}
\usepackage{setspace}
\usepackage{xcolor}

\definecolor{accent}{HTML}{2F6F73}
\definecolor{darkaccent}{HTML}{204B4D}
\definecolor{textgray}{HTML}{343C3E}
\definecolor{lightaccent}{HTML}{E9F2F1}
\hypersetup{colorlinks=true,linkcolor=darkaccent,urlcolor=darkaccent}
\setstretch{1.18}
\setlength{\parindent}{0pt}
\setlist[enumerate]{leftmargin=1.65em,itemsep=0.45em,topsep=0.3em}
\setlist[itemize]{leftmargin=1.6em,itemsep=0.2em,topsep=0.2em}
\pagestyle{fancy}
\fancyhf{}
\fancyhead[L]{\small\color{textgray}Деление многочленов и схема Горнера}
\fancyhead[R]{\small\color{textgray}11.10.26}
\fancyfoot[C]{\small\thepage}
\renewcommand{\headrulewidth}{0.3pt}
\newcommand{\R}{\mathbb{R}}
\newcommand{\lessonsection}[1]{\vspace{0.35em}\section*{\color{darkaccent}#1}\vspace{-0.75em}{\color{accent}\rule{\linewidth}{0.55pt}}\vspace{0.25em}}
\newcommand{\hint}[1]{\par\smallskip\textbf{Внимание.} #1\par}
\newcommand{\mathtask}[2]{\item \textbf{#1}\\[-0.4em]\[#2\]}

\begin{document}
\begin{titlepage}
\thispagestyle{empty}
\centering
\vspace*{28mm}
{\Large\color{darkaccent}\textbf{Чек-лист}}\\[10mm]
{\Huge\bfseries Деление многочленов}\\[4mm]
{\LARGE\bfseries и схема Горнера}\\[8mm]
{\large Уравнения, метод интервалов и замена переменной}\\[22mm]
{\large 11.10.2026}\\[27mm]
\begin{minipage}{0.79\textwidth}
\centering\large
Подбор рационального корня, деление на линейный множитель,\par
проверка остатка и переход от многочлена к произведению.
\end{minipage}
\vfill
{\large\textbf{Лёвин Артём Александрович}}\\[3mm]
{\large эксклюзивно для Анны Трапезниковой}
\vspace{16mm}
\begin{tikzpicture}[remember picture,overlay]
\draw[accent!65,line width=1.1pt]
  ([xshift=23mm,yshift=15mm]current page.south west)
  .. controls ([xshift=65mm,yshift=31mm]current page.south west)
  and ([xshift=-74mm,yshift=8mm]current page.south east)
  .. ([xshift=-23mm,yshift=18mm]current page.south east);
\end{tikzpicture}
\end{titlepage}

\lessonsection{1. Что нужно знать до деления}
Для многочлена $P(x)$ число $r$ является корнем, если $P(r)=0$. По теореме Безу:
\[
\boxed{P(r)=0\quad\Longleftrightarrow\quad P(x)=(x-r)Q(x)}.
\]
Значит, найденный корень позволяет понизить степень многочлена на единицу.

\textbf{Как подобрать корень.} Если старший коэффициент многочлена равен $1$, а все коэффициенты --- целые числа, каждый \emph{целый} корень делит свободный член. Например, для $P(x)=x^3-x^2-2x+2$ достаточно начать с $\pm1,\pm2$.
\[
P(1)=1-1-2+2=0 \quad\Longrightarrow\quad x-1\text{ --- множитель.}
\]
Для общего многочлена с целыми коэффициентами $a_nx^n+\cdots+a_0$ рациональный корень $p/q$ в несократимой дроби удовлетворяет $p\mid a_0$, $q\mid a_n$. Если $a_0=0$, сначала вынесите $x$.

\hint{Перебор делителей даёт \emph{кандидатов} на корень. Если $P(r)\ne0$, это число не является корнем.}

\lessonsection{2. Деление многочлена уголком}
Делим $P(x)=x^3-x^2-2x+2$ на $x-1$.
\begin{enumerate}
\item Выбираем первое слагаемое частного: $x^3:x=x^2$.
\item Умножаем и вычитаем:
\[
(x^3-x^2-2x+2)-x^2(x-1)=-2x+2.
\]
\item Берём следующее слагаемое: $(-2x):x=-2$.
\item Повторяем вычитание: $(-2x+2)-(-2)(x-1)=0$.
\end{enumerate}
Следовательно, частное $Q(x)=x^2-2$, остаток $0$ и
\[
\boxed{x^3-x^2-2x+2=(x-1)(x^2-2)}.
\]
Уравнение $P(x)=0$ имеет три корня:
\[
\boxed{x=1,\qquad x=-\sqrt2,\qquad x=\sqrt2.}
\]
\hint{При вычитании меняйте знаки всех слагаемых.}

\clearpage
\lessonsection{3. Схема Горнера: те же вычисления короче}
Рассмотрим $P(x)=x^3-13x-12$. Проверка показывает $P(-1)=0$, значит, делим на $x+1$. Отсутствующий член $x^2$ обязательно записываем с коэффициентом $0$.

\begin{center}
\renewcommand{\arraystretch}{1.33}
\begin{tabular}{c@{\hspace{8mm}}rrrr}
\toprule
$r=-1$ & $1$ & $0$ & $-13$ & $-12$\\
\midrule
Произведения & & $-1$ & $1$ & $12$\\
\addlinespace[2pt]
\textbf{Результат} & $\mathbf{1}$ & $\mathbf{-1}$ & $\mathbf{-12}$ & $\mathbf{0}$\\
\bottomrule
\end{tabular}
\end{center}
\textbf{Правило заполнения:} первый коэффициент сносим; каждый следующий результат умножаем на $r$ и прибавляем следующий исходный коэффициент. Последнее число --- остаток, остальные --- коэффициенты частного.
\[
\boxed{P(x)=(x+1)(x^2-x-12)=(x+1)(x-4)(x+3).}
\]
\hint{Остаток равен $P(r)$. Если заранее точно установлено, что $P(r)=0$, ненулевой остаток указывает на ошибку вычислений. Иначе он лишь означает, что $r$ не корень.}

\lessonsection{4. От разложения к методу интервалов}
На занятии рассматривалось неравенство
\[
x(x^3-13x-12)\ge0.
\]
После разложения получаем
\[
x(x+1)(x-4)(x+3)\ge0.
\]
Нули произведения: $-3$, $-1$, $0$, $4$. Все имеют нечётную кратность, поэтому знак при переходе через каждый нуль меняется.

\begin{center}
\begin{tikzpicture}[x=1.16cm,y=1cm,>=Stealth]
\draw[->,thick,accent] (-4,0)--(5.2,0) node[below right] {$x$};
\foreach \x/\v in {-3/-3,-1/-1,0/0,4/4}{
 \fill[darkaccent] (\x,0) circle (2.4pt);
 \node[below=5pt] at (\x,0) {$\v$};
}
\node[above=5pt] at (-3.55,0) {$+$};
\node[above=5pt] at (-2,0) {$-$};
\node[above=5pt] at (-0.5,0) {$+$};
\node[above=5pt] at (2,0) {$-$};
\node[above=5pt] at (4.6,0) {$+$};
\end{tikzpicture}
\end{center}
Учитывая нестрогий знак, включаем все нули:
\[
\boxed{x\in(-\infty;-3]\cup[-1;0]\cup[4;+\infty).}
\]
\hint{Знаки чередуются лишь при прохождении корня нечётной кратности. У корня чётной кратности знак произведения не меняется.}

\lessonsection{5. Повторяющееся выражение: замена переменной}
Пусть $u=x^3-x^2$. Для неравенства
\[
(x^3-x^2)^2+8(x^3-x^2)-48<0
\]
выполняем замену и получаем обычное квадратное неравенство:
\[
u^2+8u-48<0\ \Longleftrightarrow\ (u+12)(u-4)<0.
\]
Отсюда $-12<u<4$. Возвращаемся к исходной переменной:
\[
\begin{cases}
x^3-x^2>-12,\\
x^3-x^2<4.
\end{cases}
\]
Каждое неравенство раскладывается с помощью найденного корня:
\[
\begin{aligned}
x^3-x^2+12&=(x+2)(x^2-3x+6),\\
x^3-x^2-4&=(x-2)(x^2+x+2).
\end{aligned}
\]
Оба квадратных множителя положительны при всех $x\in\R$ (их дискриминанты $-15$ и $-7$). Поэтому первое условие эквивалентно $x>-2$, второе --- $x<2$. Ответ:
\[
\boxed{x\in(-2;2).}
\]
\textbf{Главное различие:} для системы берите \emph{пересечение} решений; для нескольких альтернативных случаев --- \emph{объединение}.

\lessonsection{6. Краткая самопроверка}
\begin{itemize}
\item Все степени записаны в убывающем порядке, пропущенные коэффициенты равны нулю.
\item Делитель имеет вид $x-r$ с правильным знаком.
\item Частное имеет степень на единицу меньше, остаток совпадает с $P(r)$.
\item Раскрытие скобок возвращает исходный многочлен.
\item В неравенстве учтены кратности корней и строгость знака.
\item После замены переменной решение переведено обратно к $x$.
\end{itemize}

\clearpage
\thispagestyle{fancy}
\begin{center}
{\Large\bfseries\color{darkaccent}Алгоритм: от корня к разложению}
\end{center}
\vspace{5mm}
\begin{center}
\begin{tikzpicture}[
 every node/.style={font=\small,align=center},
 startstop/.style={ellipse,draw=accent,line width=0.9pt,minimum width=43mm,minimum height=10mm,inner sep=3mm},
 act/.style={rounded corners=2mm,rectangle,draw=accent,line width=0.9pt,text width=77mm,minimum height=12mm,inner sep=3mm},
 dec/.style={diamond,aspect=3.7,draw=accent,line width=0.9pt,text width=44mm,inner sep=1.2mm},
 arrow/.style={-{Stealth[length=2.5mm]},line width=0.85pt,accent},
 node distance=11mm and 23mm]
\node[startstop] (start) {Дан многочлен $P(x)$};
\node[act,below=of start] (norm) {Упорядочьте степени; восстановите нулевые коэффициенты; вынесите общий множитель};
\node[act,below=of norm] (candidate) {Найдите кандидата $r$ на рациональный корень и вычислите $P(r)$};
\node[dec,below=of candidate] (check) {$P(r)=0$?};
\node[startstop,right=18mm of check,text width=37mm] (other) {Следующий кандидат или другой метод};
\node[act,below=13mm of check] (divide) {Разделите $P(x)$ на $x-r$ уголком или по схеме Горнера};
\node[act,below=of divide] (factor) {Запишите $P(x)=(x-r)Q(x)$; проверьте умножением};
\node[act,below=of factor] (end) {Решайте уравнение по множителям либо неравенство методом интервалов};
\node[startstop,below=of end] (finish) {Проверьте ответ};
\draw[arrow] (start)--(norm);
\draw[arrow] (norm)--(candidate);
\draw[arrow] (candidate)--(check);
\draw[arrow] (check)--node[left,pos=0.5]{да}(divide);
\draw[arrow] (check)--node[above,pos=0.5]{нет}(other);
\draw[arrow] (divide)--(factor);
\draw[arrow] (factor)--(end);
\draw[arrow] (end)--(finish);
\end{tikzpicture}
\end{center}

\clearpage
\lessonsection{Тренировочный блок \texorpdfstring{\textit{\normalsize Путь А}}{Путь А}}
\textbf{10 заданий.} Выполните вычисления самостоятельно. Там, где указан метод, используйте именно его. Рекомендуется оформлять промежуточные шаги.
\begin{enumerate}
\mathtask{Найдите один целый корень, разделите на соответствующий множитель и решите уравнение.}{x^3-x^2-4x+4=0.}
\mathtask{Выполните деление уголком на $x+1$ и решите уравнение.}{x^3+2x^2-5x-6=0.}
\mathtask{Разложите многочлен на множители и найдите все корни.}{x^3-4x^2+x+6=0.}
\mathtask{Используйте схему Горнера с $r=1$, затем решите уравнение.}{x^3-6x^2+11x-6=0.}
\mathtask{Восстановите пропущенный коэффициент и решите уравнение схемой Горнера.}{x^3-7x+6=0.}
\mathtask{Найдите рациональный корень, примените схему Горнера и решите уравнение.}{2x^3-3x^2-8x+12=0.}
\mathtask{Разложите на множители и решите неравенство методом интервалов.}{x^3-4x^2+x+6\le0.}
\mathtask{Сначала вынесите $x$, затем примените метод интервалов.}{x(x^3-4x^2-x+4)\ge0.}
\mathtask{Учтите кратность корня при расстановке знаков.}{(x-1)^2(x+2)\ge0.}
\mathtask{Сделайте замену $u=x^3-x^2$ и решите неравенство.}{(x^3-x^2)^2-2(x^3-x^2)-8<0.}
\end{enumerate}

\clearpage
\lessonsection{Ответы для самопроверки}
\renewcommand{\arraystretch}{1.42}
\begin{tabularx}{\linewidth}{@{}cX@{}}
\toprule
№ & Ответ\\
\midrule
1 & $x\in\{-2;1;2\}$\\
2 & Частное при делении на $x+1$: $x^2+x-6$; $x\in\{-3;-1;2\}$\\
3 & $x\in\{-1;2;3\}$\\
4 & Частное: $x^2-5x+6$; $x\in\{1;2;3\}$\\
5 & $x\in\{-3;1;2\}$\\
6 & $x\in\{-2;\frac32;2\}$\\
7 & $x\in(-\infty;-1]\cup[2;3]$\\
8 & $x\in(-\infty;-1]\cup[0;1]\cup[4;+\infty)$\\
9 & $x\in[-2;+\infty)$\\
10 & $x\in(-1;2)$\\
\bottomrule
\end{tabularx}

\bigskip
\textbf{Проверьте именно метод:} в заданиях 2, 4 и 6 должны появиться корректные коэффициенты частного; в заданиях 7--9 --- объяснение знаков; в задании 10 --- переход от $u$ обратно к $x$.

\end{document}
